\section{On the board}\label{sec:board}

\lsfig[t]{board}{How the LS core uses the Nexys A7-100T. Switches choose the program, the run mode, the display view and the input image; the LEDs show the FSM state, the stream valid bits and the output byte; the display shows any of seven views, here the cycle count of program S (\texttt{0x70} = 112).}

\subsection{Inputs and outputs}

The core uses four of the board's peripherals: slide switches, push buttons, LEDs and the 8-digit seven-segment display (\cref{fig:board}). Pins come from Digilent's master constraints file for the board, which must be credited in the report; the table is for cross-checking.

\begin{center}\small
\begin{tabularx}{\linewidth}{@{}l l L@{}}
\toprule
Port & Pin(s) & Use \\
\midrule
\sig{clk} & E3 & 100\,MHz, the only clock; \texttt{create\_clock} with a 10\,ns period \\
BTNU & M18 & synchronous reset, through a two-flip-flop synchroniser \\
BTNC & N17 & go / step: synchroniser, 10\,ms debounce, one-clock pulse \\
SW15..13 & V10, U11, U12 & program select, read at reset \\
SW12..11 & H6, T13 & run mode \\
SW10..8 & R16, U8, T8 & display view (SW9 and SW8 are 1.8\,V inputs: \texttt{LVCMOS18}) \\
SW7..0 & R13, U18, T18, R17, R15, M13, L16, J15 & \insn{IN} data; SW7..6 choose the image \\
\bottomrule
\end{tabularx}
\end{center}

\paragraph{LEDs} LED15..12 light in \stname{FETCH}, \stname{DECODE}, \stname{EXECUTE}, \stname{WRITEBACK}; LED11 in \stname{STREAM}; LED10 in either drain state; LED9 and LED8 show \sig{V1} and \sig{V2}; LED7..0 show \sig{OUTR}, the neuron's output byte.

\paragraph{Display} Eight common-anode digits share one set of segment lines. The scanner lights one digit at a time for 1\,ms (a tick every 100,000 clocks), so the whole display refreshes every 8\,ms, faster than the eye can see; segments and anodes are active-low, and a short blank between digits prevents ghosting. Until the 8-digit driver of Lab 11 exists, one digit or the LEDs are enough.

\begin{center}\small
\begin{tabular}{@{}l l@{\qquad}l l@{}}
\toprule
\multicolumn{2}{@{}l}{View (SW10..8)} & \multicolumn{2}{l@{}}{Program (SW15..13)} \\
\midrule
0 & \sig{PC}, state, \sig{IR} & \texttt{000} & S64 \\
1 & \sig{R0}--\sig{R3} & \texttt{001} & U64 \\
2 & \sig{ACC}, \sig{OUTR} & \texttt{010} & L64 \\
3 & \sig{CNT}, \sig{P0}, \sig{P1} & \texttt{011} & B64 \\
4 & cycle counter (8 hex digits) & \texttt{100} & S32 \\
5 & instruction counter & \texttt{101} & U32 \\
6 & \sig{PROD}, \sig{V1}, \sig{V2}, \sig{Z}, \sig{N} & \texttt{110} & L32 \\
 & & \texttt{111} & \insn{MAC \#8} demo \\
\bottomrule
\end{tabular}
\end{center}

\subsection{The demonstration, in three minutes}

\begin{enumerate}[leftmargin=1.6em]
\item \textbf{Watch a stream happen.} Cycle-step the \insn{MAC \#8} program with views 3 and 2. LED15..12 show \stname{FETCH}, \stname{DECODE}, then LED11 lights for \stname{STREAM}; with each press \sig{CNT} falls and \sig{P0}, \sig{P1} rise by one; \sig{V1} lights one press into the stream and \sig{V2} one press later; two drain presses finish it. The board matches \cref{fig:timing-mac4} press by press.
\item \textbf{Race the four programs.} Run mode, view 4, reset between runs: the display shows the counts of S64, U64, L64 and B64, and LED7..0 show the same output byte every time.
\item \textbf{Change the data, not the time.} Flip SW7..6 to another image: the output byte changes, the cycle counts do not.
\item \textbf{The slope.} S32 and L32: thirty-two fewer inputs save 32 cycles with the stream and 256 with the loop.
\end{enumerate}

\subsection{Demo data}

A hand-made neuron works better on the board than a trained one: a template detector for one digit, with small positive weights on the digit's strokes and small negative weights around them, drawn on squared paper. Four stored images give clearly different sums. The shift $sh$ is the smallest that brings the largest sum within 127; the bias and ReLU are set so the four output bytes differ and none is \texttt{00} or \texttt{7F} (a saturated value would hide errors). All four outputs are written on the prediction sheet before the board is switched on. A negative check for the requantiser: $\mathrm{ACC} = -301$, $sh = 2$ gives $-76$ (\texttt{B4}) with ReLU off and \texttt{00} with ReLU on.
